\answer{ex:22:1}{(a)~$32\%$와 $100\%$. (b)~$R_\uparrow-R_\downarrow\ge2\sqrt{2000\,\text{Hz}/1\,\text{s}}\approx89$~Hz, 전체 빈도의 $4.5\%$에 불과하다.}
\answer{ex:22:2}{(a)~$n\ge4(p_1q_1+p_2q_2)/\Delta^2\approx560$. (b)~$\approx56$.}
\answer{ex:22:3}{(a)~세부 균형: $\rho PK$가 대칭이다. (b)~미스 비율은 조화평균 $1/\langle1/P\rangle=0.18$로, 되먹임이 없을 때의 $0.5$와 대비된다.}
\answer{ex:22:4}{(a)~넓이 $4h^2=0.04$인 평행사변형($p$의 제한을 고려하면 수치로 $0.045$). (b)~$\approx0.18$.}
\answer{ex:22:5}{$0.49$: 일부 배양은 거의 언제나 미스인 영역으로 빠져나가 거기서 떠돈다. 제한을 두면 $1.00$: 유계 집합 위에서는 밀도 $\propto1/P$가 흡수 영역에 모인다.}
\answer{ex:22:6}{(a)~수준이 $\ge5$개 필요하다. 전극 $8$개면 충분하다. (b)~아니다. $r_{\max}\ge25/T=100$~Hz가 필요하다.}
\answer{ex:22:7}{열린 문제. 무작위 탐색 시간은 $d$에 따라 대략 부피 비 $(L/h)^d$처럼 늘고, 오차 부호를 쓰는 탐색은 $d$에 선형으로 는다. 두 가설을 가르는 것은 뒤바꾸기 실험이다. 미스 뒤에는 예측 가능하지만 강한 자극을, 히트 뒤에는 예측 불가능하지만 약한 자극을 준다. 그룹마다 배양이 수십 개 필요하다.}
